\section{引言}
\sectionclaim{完整证据覆盖而非主题相关性，决定了加速器复合问题能否得到有据回答。}

加速器工程查询会把设备、动作、条件和版本压缩进同一句问题。SHINE 主加速器的“工作模式与最终能量”要求两个事实；低温设备保护要求区分主联锁 PLC 与 CPLD；IOC 启动信息要求完整字段清单。这些问题的共同特征是：检索结果可以与问题高度相关，同时遗漏其中一个必需条件。以可执行的证据完整性约束连接检索与生成，是工程知识服务的核心问题。

现有系统在我们的语料中暴露出三种具体失效模式。第一，BGE-small-zh 直接编码时有 376/4,532 个输入超过 512 tokens，字符分块与模型预算发生错配。第二，主加速器父块完整包含所需证据，但实际送入重排的片段只保留局部内容。第三，初始摘录式回答给出能量却漏掉模式，仍返回完整状态。三种失效分别发生在输入表示、候选评分和回答覆盖层，单一相似度分数难以定位责任。

我们提出以证据契约贯通这三个层次的系统主张：来源身份固定，模型表示可变，完整性逐层可计算。系统将规范父块作为引用单位，将模型专用片段作为编码单位；以文档、物理页和原文摘录约束证据命中；再把用户要求、选中片段、生成主张与引用建立显式绑定。本文只覆盖静态文档问答。

领域语料使这种设计具有实际价值。实验物理与工业控制系统（EPICS）通过输入输出控制器（IOC）和过程变量（PV）组织设备访问，其技术标识需要精确保留\citep{epics,epicsdocs}；SHINE 控制架构\citep{shinecontrol}与 European XFEL 设计资料\citep{xfel}又包含不同装置、年代和参数定义。系统必须在中英文术语对应的同时恢复事实的所属对象与文档版本。

本文给出三条贡献：
\begin{enumerate}
\item 提出引用与编码解耦的证据契约，在 13 份 PDF、2,178 页和 4,532 个可检索父块上实现稳定来源与 token 预算约束。
\item 比较 11 种检索配置，以逐问题分析定位异构候选互补，将完整 Hit@5 由 18/21 提高到 20/21，净增 2 条且无反向翻转。
\item 实现需求到原文的可追溯问答链，完成 32/32 条请求验证，平均 2.51 s，并定位 20/21 父块覆盖与 19/21 片段覆盖的粒度差异。
\end{enumerate}

实验结果把系统选择建立在明确的成本与收益上。多查询改写使 E5 完整命中由 18 条降到 15 条，每条查询的平均检索调用量升至 2.844 倍，最终系统据此保留原始查询。双模型重排的两条新增命中来自 Qwen3 补充的候选，而 E5 原有 18 条命中全部保留。

逐项追溯还把回答完整性变成可诊断对象。冻结字段清单显示，一条 IOC 答案在内部核查中被标为支持，却遗漏部分必需字段。这个实例揭示了“来源支持某句话”与“回答覆盖用户全部要求”的判定差异，推动系统评价同时保留检索、片段与回答三个层次的证据。
